\section{Introduction to the Topic and Historical Context}

\subsection{Introduction}

In November 1975, Latin America stands at the center of one of the Cold War's most volatile fronts.  Across the continent, military governments, revolutionary movements,  intelligence agencies and foreign powers  are disputing one another in a “battle” whose outcome will determine the political future of the Western Hemisphere.

This instability responds to deep systematic friction. Following a series of military coups in the Southern Cone,  traditional policing has transformed into a coordinated architecture of state terror. Driven by the National Security Doctrine, governments have redefined the internal enemy to include left-wing guerrilla movements and political dissidents. So by sharing clandestine databases and specialized execution squads they ensure that crossing a border no longer offers safety.

Said National Security Doctrine emerged as the ideological backbone of Latin American military dictatorships in the 70s,  by reshaping governance into a military crusade against the  perceived internal enemies. This Doctrine is born from the Cold War of paranoia,  redefining democracy as a fragile structure vulnerable to communist infiltration. Left-wing movements, student activists, labor unions and even moderated reformists are  perceived as existential threats justifying the violence seizure of power by military juntas.  Again, this Doctrine provides the moral and strategic framework for regimes to dismantle democratic institutions replacing them with systems designed to wage total war against its own citizens.

In 1959, the Cuban Revolution brought the Cold War directly to the Western Hemisphere, prompting an immediate response from the United States with the Alliance for Progress in 1961, a development and counterinsurgency program designed to prevent other revolutions from following the same path. In 1964, the military coup in Brazil inaugurated a new era of dictatorships under the National Security Doctrine, and by 1973, the coups in Chile and Uruguay consolidated a bloc of authoritarian regimes across the Southern Cone.

For the United States, this conflict's  importance goes far beyond territory. Every time an insurgency successfully takes place, the US faces extreme in house criticism. As concern regarding Soviet and Cuban influence grows larger each day, this dispute shifts perspective. For the US, this conflict is not about purely stopping communism anymore, but also about stating their power and how much they are willing to put on the table.

\subsection{Key Terms}

\begin{guidetable}[Y{0.46}Y{1.54}]{2}
  \textbf{National Security Doctrine} & The ideological backbone, born from Cold War paranoia, that has reshaped Latin American governance into a military crusade, redefining democracy as a fragile structure vulnerable to communist infiltration and casting student activists, labor unions, and political dissidents as existential threats. \\
  \midrule
  \textbf{Southern Cone} & Southern region of South America: Chile, Argentina, Uruguay, Paraguay, and extended partners. Where a series of military coups have consolidated a bloc of right-wing authoritarian regimes locked in a battle to determine the political future of the Western Hemisphere. \\
  \midrule
  \textbf{Central Intelligence Agency (CIA)} & The primary independent intelligence apparatus of the United States, established under the National Security Act of 1947, tasked with evaluating, correlating, and navigating the delicate boundary between gathering vital information and managing foreign interventions. \\
  \midrule
  \textbf{Intelligence Collection vs Covert Action} & Intelligence collection serves as the agency's lawful foundation of gathering information, running informants, and monitoring communications, whereas covert action represents clandestine operations designed to covertly manipulate foreign political and military conditions, legally demanding explicit presidential authorization. \\
  \midrule
  \textbf{Presidential Finding} & The indispensable legal directive issued directly by the President of the United States, serving as the sole formal mandate required before the Agency can unleash covert operations beyond American borders. \\
  \midrule
  \textbf{Liaison relationships} & Clandestine, bilateral intelligence channels connecting the Agency to individual foreign security services, through which technical equipment, tactical training, and sensitive information flow across the hemisphere. \\
  \midrule
  \textbf{Church committees} & The Congressional inquiry operating in Washington to expose decades of secretive U.S. interventions, covert funding, and illegal foreign operations, placing the Agency under direct political strain and heightened fear of exposure. \\
  \midrule
  \textbf{DINA \& SIDE} & The secret police agencies of Chile and Argentina, headed by figures like Manuel Contreras, that serve as the primary execution squads and principal architects driving the machinery of transnational repression. \\
  \midrule
  \textbf{School of the Americas} & The U.S.-run military training facility where foreign officers are instructed in counterinsurgency tactics and doctrine, effectively shaping the operational vocabulary and coercive methods to be deployed in the battle against internal enemies. \\
\end{guidetable}

\subsection{The CIA: history and competences}

The Central Intelligence Agency was created by the National Security Act of 1947, which established it as an independent agency operating under the direction of the National Security Council. Its original statutory purpose, as defined in that Act, was to "correlate and evaluate intelligence relating to national security, and provide for the appropriate dissemination of such intelligence within the Government" (National Security Archive, "The National Security Act Turns 75"). The Act was written in direct response to the intelligence failures that preceded Pearl Harbor, and it aimed to centralize what analysts call "national" intelligence, information of interest to multiple departments, while still permitting the military services, the State Department, and other agencies to keep their own "departmental" intelligence functions (Federation of American Scientists, irp.fas.org/offdocs/int010.html).

The Agency is led by the Director of Central Intelligence (DCI), who by statute is the head of the CIA and who, in practice, also serves as the President's principal intelligence adviser. Because the position depends so heavily on the President's personal trust, each administration has traditionally been given wide latitude to select its own Director. Institutionally, the CIA is organized around four main directorates: Operations (responsible for human intelligence collection and for covert action), Intelligence (responsible for analysis), Science and Technology (responsible for technical collection), and Administration.

For the purposes of this committee, the most important distinction delegates should understand is the one between intelligence collection and covert action. Intelligence collection, gathering information, running informants, monitoring communications, is understood as the Agency's core, lawful function. Covert action, by contrast, is defined as an activity intended to influence political, economic, or military conditions abroad in a way where the role of the United States government is not intended to be apparent. Covert action is not an inherent power of the CIA; it can only be undertaken "as may be approved and directed by the President". In other words, the CIA does not independently decide to overthrow a government, arm an insurgency, or run an assassination program, legally, such actions require a presidential finding, however informal that finding may be.

It is worth stressing to the committee that in November 1975, this legal architecture is under direct strain. The Church Committee, chaired by Senator Frank Church, is in the process of publicly investigating the CIA's history of covert operations, and its findings — some of which had already been released by December 1975, describe "covert United States involvement in Chile in the decade between 1963 and 1973" as "extensive and continuous," including CIA efforts to prevent Salvador Allende's election and, later, to destabilize his government (National Security Archive, "Covert Action in Chile"). The Committee also documented that President Nixon personally ordered CIA Director Richard Helms, on September 15, 1970, to "prevent Allende's inauguration by fomenting a preemptive coup," and that the Agency subsequently paid \$35,000 in what the Committee's report called "hush money" connected to the assassins of General René Schneider (National Security Archive, "Covert Action in Chile"). It should be understood that this history is precisely why the Agency, in November 1975, is acutely conscious of exposure, both political and legal, for any perceived complicity in the repression unfolding in the Southern Cone.

\subsection{Cold war in Latin America}

The Cold War is the crash of two main ideologies and therefore global superpowers that reshaped global dynamics, and Latin America stands as one of the most critical theaters of this conflict. Said main ideologies are represented by the United States as democracy, free market economies and anti-comunism, and the Soviet Union as socialism revolution, state controlled economies and anti-imperialism.

For the U.S the stakes in Latin America are crucial, losing this region to communism would mean the expansion of the Soviet influence just miles away from its borders,  threatening national security.

The Cuban revolution in 1959 served as a wake up call for the U.S. With Fidel Castro in power, Cuba aligned with the Soviet Union ideals, and by 1960, Cuba had totally embraced communism. Soviet Premier Nikita Khrushchev publicly declared that Latin America was “ ripe for revolution” signaling Moscow's interest to exploit the national instability of the continent.

The U.S had to respond and did so in two main ways: diplomatic and economic engagement, with initiatives like the Alliance for Progress in 1961, aimed to counter communist propaganda by fighting poverty and inequality, and  by strategic interventions in which the CIA has played a central role in supporting coups and training military forces in counterinsurgency tactics to suppress leftist movements.

By the beginning of the 70s, the region was deeply polarized. Some nations like Chile, under Salvador Allende, pursued socialist paths through democratic means, while others like Brazil, turned into military regimes. The Cold War in Latin America is the driving force behind the rise of authoritarianism, and the consequences that it faces now.

\subsection{The rise of the National Security Doctrine}

The National Security Doctrine emerged as the ideological foundation for the military seizure of power in Latin America. Born from the fear of internal subversion and the perceived threat of communist expansion, this doctrine redefined governance as a military crusade against the perceived enemies. It argued that democracy was way too fragile to withstand the Marxist threat and just defined violence overthrow of elected governments in the name of stability and security. This Doctrine views the internal enemy as any kind of left-wing movements, labor unions, student activists and even moderate reformists.  These collectives are cast as existential threats to the state, and a problem to take care of.

Furthermore, with the establishment of the National Security Doctrine, the states have turned to a militarized political system. The Armed Forces assume the role of guardians of national identity positioning themselves as the ultimate arbiters of political legitimacy. In this climax, repression is justified; and torture, forced disappearances and extrajudicial killings are framed as necessary tools to preserve civilization.

By the early 70s, military dictatorships have taken root in Chile and Uruguay  in 1973, and in Brazil in 1964. These regimes do not act in isolation, they share intelligence, coordinate repression and lay the ground for a transnational coalition.

\subsection{U.S strategic interest in the Southern Cone}

\subsubsection{Anti-communism as a U.S priority}

For Washington, the Cold War is a fight for every region in Latin America. The Southern Cone, though geographically distant from the United States, is strategically vital as it contains shipping routes, resource rich economies and governments that could either serve as allies or become vulnerabilities. The core U.S priority is simple, prevent any communist governments from taking root in the hemisphere. After Cuba in 1959, the fear lies in a possible domino effect: if one country falls, others may follow. This logic drives America's policy from the Alliance of Progress to the training of military officers at the School of the Americas, and gives the National Security Doctrine a powerful patron.

Since the Cuban Revolution, Havana has become a staging ground for Soviet influence in the region, offering training, funding and ideological support to the leftist movement in the continent, therefore, containing Cuba's influence is as urgent as containing the Soviets themselves. The Southern Cone is seen as particularly vulnerable; Chile, under Salvador Allende, showed that Marxism could come to power through democratic elections, and Uruguay's Tupamaros demonstrated the power of urban guerrilla warfare. Washington fears that the Southern Cone could become a second front where revolutionary movements inspired by Havana and backed by Moscow could destabilize the entire region. Every coup, every intelligence sharing agreement, and every covert operation is justified in the name of stopping this spread.

\subsubsection{Protection of U.S credibility in the Cold War}

The United States is locked in a global ideological struggle, and its credibility depends on appearing strong and reliable. If a country in its own hemisphere turned communist, it will be read as a failure of American power worldwide. This logic drove the Kennedy and Johnson administrations to deepen involvement in the region, and it is the same logic that led Nixon and Kissinger to support the coup in Chile. The Southern Cone has become a stage where the U.S must prove that it can protect its sphere of communist influence.

This creates a fundamental tension in U.S policy. The Brazilian regime of 1964, the Chilean junta of 1973, and the Uruguayan dictatorship all receive varying degrees of American support through military aid, intelligence sharing, diplomatic backing, and training. For U.S. policymakers, these regimes are not ideal partners, but they are reliable ones. They crush leftist movements, maintain order, and keep their economies open to American investment.
